\chapter{Hamiltonian fields and continuous coefficients}
\label{mht:ch:hamiltonian-coefficients}

A Hamiltonian function can serve as a field on a symplectic surface
or as an independent coefficient of another field. These uses have
different differentials and support conditions. Their formal
coordinates alone do not identify the resulting theories.

Finite Hamiltonian jets provide explicit coefficient algebras.
Their duals, quantizations and continuous cochains require separate
constructions. The first issue is already visible before
quantization: discarding a Hamiltonian jet reverses direction on its
dual. For cochains, the output jet and the input jet must be allowed
to differ.

\section{Pointed Hamiltonians and their finite quotients}

Work over $\mathbb C$ and put $\mathfrak m=(u,v)$ in
$\mathbb C[[u,v]]$. Give this ring its $\mathfrak m$-adic topology,
with discrete finite quotients. The Poisson bracket is
\begin{equation}\label{mht:eq:ham-poisson}
 \{f,g\}=\partial_uf\,\partial_vg-\partial_vf\,\partial_ug.
\end{equation}
The Hamiltonian vector field is $X_f=\{f,-\}$.
Constants have zero vector field. Thus
$\mathfrak h_{\mathrm{aff}}=\mathbb C[[u,v]]/\mathbb C$
contains every formal Hamiltonian vector field on the plane.

The pointed algebra is the Lie subalgebra
$\mathfrak h=\mathfrak m^2$: its vector fields vanish at the origin.
The filtration $F^a\mathfrak h=\mathfrak m^{a+2}$, for $a\geq0$,
satisfies
\[
 \{F^a\mathfrak h,F^b\mathfrak h\}\subset F^{a+b}\mathfrak h.
\]
For $d\geq2$, write $W_r$ for homogeneous polynomials of degree $r$
and set
\[
 \mathfrak h_d=\mathfrak m^2/\mathfrak m^{d+1}
              =\bigoplus_{r=2}^dW_r.
\]
These coefficient Lie algebras have cohomological degree zero.
Polynomial degree is a separate grading.

\begin{proposition}\label{mht:prop:pointed-jets}
The bracket descends to every $\mathfrak h_d$, and the degree
projections are Lie maps. Moreover,
\[
 \mathfrak h=\varprojlim_d\mathfrak h_d,\qquad
 \dim\mathfrak h_d=\frac{d(d+3)}2-2.
\]
Every $\mathfrak h_d$ is unimodular:
$\operatorname{Tr}_{\mathfrak h_d}(\operatorname{ad}_x)=0$
for every $x$.
\end{proposition}

\begin{proof}
The bracket sends $\mathfrak m^r\times\mathfrak m^s$ into
$\mathfrak m^{r+s-2}$. Hence $\mathfrak m^{d+1}$ is an ideal
in $\mathfrak m^2$. Compatible jets determine one formal series,
and summing $\dim W_r=r+1$ gives the dimension.

Set $E=u^2/2$, $F=-v^2/2$ and $H=-uv$. Direct differentiation gives
\[
 \{H,E\}=2E,\qquad \{H,F\}=-2F,\qquad \{E,F\}=H.
\]
On $u^{r-j}v^j$, the $H$-weights are $r-2j$.
They sum to zero. The operators $E,F$ move between distinct weights,
so their traces vanish.

Higher-degree Hamiltonians strictly raise
polynomial degree and have triangular adjoint matrices with zero
diagonal. Linearity proves unimodularity.
\end{proof}

The quadratic algebra is $\mathfrak{sl}_2$. Its presence means that
$\mathfrak h$ is not pronilpotent. The ideal
$\mathfrak m^3=F^1\mathfrak h$ is pronilpotent, since an iterated
bracket of $r$ of its elements lies in $F^r\mathfrak h$.

The first nonlinear layers exhibit this filtration. Put
$v_i=u^{3-i}v^i$, for $0\leq i\leq3$, and
$w_j=u^{4-j}v^j$, for $0\leq j\leq4$. Then
\begin{align*}
 \{H,v_i\}&=(3-2i)v_i,&
 \{E,v_i\}&=iv_{i-1},&
 \{F,v_i\}&=(3-i)v_{i+1},\\
 \{v_i,v_j\}&=3(j-i)w_{i+j-1}.
\end{align*}
Terms outside the indicated index ranges vanish.

Consequently $\mathfrak h_3=\mathfrak{sl}_2\ltimes W_3$ with
abelian $W_3$. In $\mathfrak h_4$, the cubic brackets generate
$W_4$, while $\{W_3,W_4\}=\{W_4,W_4\}=0$.
For example, $\{u^3,v^3\}=9u^2v^2$ and
$\{u^2v,uv^2\}=3u^2v^2$.

Translations change the permitted quotients:
\[
 \{u,v^{d+1}\}=(d+1)v^d.
\]
The right side survives the degree-$d$ projection although
$v^{d+1}$ does not. The same jet ideal therefore cannot be used
after adjoining the linear Hamiltonians.

\section{Residue duals and cotangent degrees}

A continuous functional from $\mathfrak h$ to discrete $\mathbb C$
annihilates some $\mathfrak m^{d+1}$. Its continuous dual is
\[
 P=\mathfrak h^\vee_{\mathrm{cont}}
    =\varinjlim_d\mathfrak h_d^\vee
    =\bigoplus_{a+b\geq2}\mathbb C\rho_{ab}.
\]
Here $\rho_{ab}$ extracts the coefficient of $u^av^b$.
It can be represented by the principal part
$u^{-a-1}v^{-b-1}\,du\,dv$, with residue convention
\[
 \operatorname{Res}_0(u^iv^j\rho_{ab})=\delta_{ia}\delta_{jb}.
\]
The algebraic direct sum has the discrete linear topology.
The residue pairing with the formal product is continuous.

Use the coadjoint convention
$\langle\operatorname{ad}_f^*\rho,g\rangle
=-\langle\rho,\{f,g\}\rangle$. The monomial bracket gives
\begin{equation}\label{mht:eq:ham-coadjoint}
 \operatorname{ad}_{u^iv^j}^*\rho_{ab}
  =\bigl(j(a+1)-i(b+1)\bigr)
       \rho_{a-i+1,b-j+1}.
\end{equation}
The expression is zero when an index is negative or the output
degree is less than two. Indeed, pairing with $u^kv^\ell$ forces
$k=a-i+1$ and $\ell=b-j+1$, and negating the bracket coefficient
gives the displayed factor.

\begin{example}\label{mht:ex:cotangent-jet-projection}
The simultaneous degree projection
\[
 \mathfrak h_4\ltimes\mathfrak h_4^\vee[-1]
 \longrightarrow
 \mathfrak h_3\ltimes\mathfrak h_3^\vee[-1]
\]
is not a Lie map. Formula \eqref{mht:eq:ham-coadjoint} gives
$\operatorname{ad}_{u^3}^*\rho_{22}=-9\rho_{03}$.
The projection kills $\rho_{22}$ while preserving $u^3$ and
$\rho_{03}$. Thus projecting their bracket differs from bracketing
their projections.
\end{example}

The actual dual of $\mathfrak h_4\to\mathfrak h_3$ is the injection
$\mathfrak h_3^\vee\to\mathfrak h_4^\vee$. One may retain both
systems, with their opposite directions. They do not form the
simultaneous inverse system in the example.

This distinction also controls completion. For a countable basis
$(e_i)$, the vectors $e_i$ tend to zero in the coordinate product.
A continuous bilinear pairing of two such products cannot satisfy
$B(e_i,e_j)=\delta_{ij}$, since $B(e_i,e_i)=1$.
Reciprocal rescaling of two dual bases leaves their contractions
unchanged. In particular, contracting the first $n$ diagonal
terms still gives $n$. Weights alone do not define that infinite
contraction.

For a finite-dimensional Lie algebra $\mathfrak l$, two cotangent
shifts have different Chevalley degrees. With
$V[k]^j=V^{j+k}$ and
$C^\bullet_{\mathrm{CE}}(\mathfrak g)
=\operatorname{Sym}(\mathfrak g^\vee[-1])$, put
\[
 \mathfrak g_-=\mathfrak l\ltimes\mathfrak l^\vee[-1],
 \qquad
 \mathfrak g_+=\mathfrak l\ltimes\mathfrak l^\vee[1].
\]
The shifted dual is abelian, with its coadjoint action.
The exterior ghost $c^a$ has degree one in both cases.
The coordinate dual to the cotangent summand has degree zero for
$\mathfrak g_-$ and degree two for $\mathfrak g_+$.

For the former there is an isomorphism of commutative complexes
\[
 C^\bullet_{\mathrm{CE}}(\mathfrak g_-)
   =C^\bullet_{\mathrm{CE}}
      (\mathfrak l,\operatorname{Sym}\mathfrak l).
\]
Its differential is the adjoint Chevalley differential.

Subtracting twice the cotangent weight from the degree in
$C^\bullet_{\mathrm{CE}}(\mathfrak g_+)$ gives the same coordinate
degrees. This is an explicit regrading. It is not an identification
of the two original graded field theories.

There is also a finite BV model whose antifields have degree $-2$.
Choose $[e_a,e_b]=f^c{}_{ab}e_c$, ghosts $c^a$ of degree one and
antifields $b_a$ of degree $-2$. With left odd derivatives, set
\[
 \Delta=\sum_a\partial_{b_a}\partial_{c^a},\qquad
 S=-\frac12\sum_{a,b,c}b_a f^a{}_{bc}c^bc^c.
\]
The BV bracket is fixed by the second-order identity
\[
 \Delta(FG)=(\Delta F)G+(-1)^{|F|}F\Delta G
                         +(-1)^{|F|}\{F,G\}.
\]

\begin{proposition}\label{mht:prop:finite-hamiltonian-bv}
If $\mathfrak l$ is unimodular, then
$\{S,S\}/2+\hbar\Delta S=0$. In particular this holds for every
$\mathfrak h_d$ with its constant Berezinian.
\end{proposition}

\begin{proof}
The Hamiltonian differential sends
$c^a$ to $-f^a{}_{bc}c^bc^c/2$ and sends the $b$ variables by
the cotangent lift. Jacobi gives its square zero.
Equivalently, the coefficients of $\{S,S\}$ are the Jacobi sums,
so they vanish. Direct contraction gives
\[
 \Delta S=\sum_k\operatorname{Tr}(\operatorname{ad}_{e_k})c^k.
\]
Unimodularity annihilates this term.
\end{proof}

The finite odd cotangent model has its declared coordinates,
measure and differential. Its quantum master equation supplies no
measure or diagonal contraction on an unrestricted infinite
coefficient product.

\section{Horizontal physical jets}

Let $z=(z_1,z_2)$ be holomorphic coordinates on a flat chart $U$.
Introduce independent formal variables $w_1,w_2$ and odd variables
$\vartheta_1,\vartheta_2$ of degree one. Define
\[
 \mathcal J=\mathcal O(U)[[w_1,w_2]]
               \otimes\Lambda(\vartheta_1,\vartheta_2),\qquad
 \nabla=\sum_i\vartheta_i(\partial_{z_i}-\partial_{w_i}).
\]
The formal coefficients have the product topology, using compact-set
holomorphic seminorms on each coefficient. Exterior degree is
auxiliary and adds no dimension to the physical integration measure.

Define $T=\exp(w\cdot\partial_z)$ coefficientwise, and let
$j(f)=Tf$ be the formal Taylor series of $f(z+w)$.
Let $p(F)$ evaluate the exterior-degree-zero part at $w=0$.
Put $\delta=\sum_i\vartheta_i\partial_{w_i}$ and
$\kappa=\sum_iw_i\iota_{\vartheta_i}$, where $\iota$ is the left
exterior derivative. On $w$-degree $m$ and exterior degree $r$, set
\[
 h_0=-\frac{\kappa}{m+r}\quad(m+r>0),\qquad
 h_0=0\quad(m+r=0),\qquad h=Th_0T^{-1}.
\]

\begin{proposition}\label{mht:prop:physical-spencer}
These continuous maps satisfy
\[
 pj=1,\qquad \nabla h+h\nabla=1-jp,\qquad
 h^2=hj=ph=0.
\]
The Taylor inclusion preserves the constant Poisson bracket and
its formal Moyal product.
\end{proposition}

\begin{proof}
The anticommutator $\delta\kappa+\kappa\delta$ is multiplication
by $m+r$. Thus $h_0$ contracts $-\delta$ onto the constants in
$w,\vartheta$. Differentiation gives
$\nabla=T(-\delta)T^{-1}$, and evaluation at $w=0$ commutes with
$T^{-1}$. Conjugating proves all the contraction identities.

At fixed output $w$-degree, $T$ and $T^{-1}$ use finitely many
coefficient derivatives. The operator $h_0$ raises that degree by
one and lowers exterior degree by one. These facts prove continuity.
Finally $\partial_{w_i}j(f)=j(\partial_{z_i}f)$.
Applying this equality to each derivative proves the Poisson and
Moyal assertions coefficientwise.
\end{proof}

For example,
$j(z_1^2)=z_1^2+2z_1w_1+w_1^2$.
Evaluation on arbitrary independent coefficients is not Poisson:
\[
 p\{w_1,w_2\}_w=1,\qquad \{p(w_1),p(w_2)\}_z=0.
\]
The contraction therefore compares geometric functions with the
horizontal complex. Imposing that complex changes a theory of
unconstrained independent coefficient fields. The same formulas
work for smooth coefficients and the total mixed differential,
with its tensor signs and coefficientwise support.

\section{The abstract Rees algebra and the Moyal product}

For a finite-dimensional Lie algebra $\mathfrak l$, let
$R_N=\mathbb C[\hbar]/(\hbar^{N+1})$ and define
\[
 U_N(\mathfrak l)=
 T_{R_N}(R_N\otimes\mathfrak l)/
       (xy-yx-\hbar[x,y]).
\]
Here $N$ is parameter precision, unrelated to matrix rank.
The Poincar\'e--Birkhoff--Witt map is the module isomorphism
\[
 \operatorname{sym}_N:
 \operatorname{Sym}(\mathfrak l)\otimes R_N
 \longrightarrow U_N(\mathfrak l),\qquad
 x_1\cdots x_r\longmapsto
 \frac1{r!}\sum_{\sigma\in S_r}x_{\sigma(1)}\cdots x_{\sigma(r)}.
\]
The source here has its ordinary symmetric product. Symmetrization
is not generally an algebra map.

For completeness, choose an ordered basis of $\mathfrak l$.
Interchanging an inverted neighboring pair adds its shorter
$\hbar$-bracket term. Reduction terminates by word length and then
lexicographic order at fixed length. The only overlapping
ambiguity has three factors, and its difference reduces to the
Jacobi sum. Disjoint interchanges commute.
Thus ordered monomials form a basis, over $R_N$ as well as over
$\mathbb C[\hbar]$.

PBW degree differs from powers of the augmentation ideal.
For example, let $\mathfrak n$ be the Heisenberg Lie algebra
with basis $x,y,z$, bracket $[x,y]=z$ and central $z$.
In its ordinary enveloping algebra, the augmentation sends these
three generators to zero. Its kernel $I$ contains
$z=xy-yx\in I^2$, although $z$ has PBW degree one.
The two filtrations therefore define different truncations.

The coefficient limit is
$U_\hbar(\mathfrak l)=\varprojlim_NU_N(\mathfrak l)$.
Its PBW module is $(\operatorname{Sym}\mathfrak l)[[\hbar]]$:
each coefficient is polynomial, but its degree can increase with
parameter order. For example, $\sum_{n\geq0}\hbar^n x^n$
belongs to this module and need not belong to
$\mathbb C[[\hbar]][x]$.

Quantizing functions on the symplectic plane gives a different
operation. On $\mathbb C[u,v][[\hbar]]$, define
\begin{equation}\label{mht:eq:ham-moyal}
 f\star g=
 \sum_{r\geq0}\frac{\hbar^r}{2^rr!}
 \sum_{j=0}^r(-1)^j\binom rj
  (\partial_u^{r-j}\partial_v^jf)
  (\partial_u^j\partial_v^{r-j}g).
\end{equation}
For polynomial inputs the sum terminates. In formal $\hbar$
series each coefficient still has finitely many contributions.
The commuting constant bidifferential operators on three factors
prove associativity by multiplying their exponentials.

\begin{proposition}\label{mht:prop:pointed-moyal-failure}
Quadratic Hamiltonians have
$[f,g]_\star=\hbar\{f,g\}$. The full pointed polynomial subspace
is not closed under $\hbar^{-1}[-,-]_\star$, even modulo constants:
\begin{align*}
 [u^3,v^3]_\star&=9\hbar u^2v^2+\frac32\hbar^3,\\
 [u^3,v^4]_\star&=12\hbar u^2v^3+6\hbar^3v.
\end{align*}
\end{proposition}

\begin{proof}
Reversing the inputs multiplies the derivative-order-$r$ term by
$(-1)^r$. The commutator retains precisely the odd orders.
Orders at least three vanish on quadratics. For $u^a,v^b$ the
surviving coefficient is
\[
 2\frac{\hbar^r}{2^rr!}
   \frac{a!}{(a-r)!}\frac{b!}{(b-r)!}\,u^{a-r}v^{b-r}.
\]
Taking $(a,b)=(3,3)$ and $(3,4)$ gives the formulas.
The second normalized commutator retains the linear term
$6\hbar^2v$ modulo constants.
\end{proof}

The Rees algebra of $\mathfrak h_d$ uses its specified truncated
classical bracket. It does not make the pointed subspace of
Moyal functions closed. A quantum Hamiltonian action retaining
all polynomial functions belongs to the larger affine recipient.

\section{The finite relative cochain comparison}

Put $A=U_N(\mathfrak l)$ and let $M$ be an $A$-bimodule.
The relative Hochschild cochains are
\[
 C^n_{\mathrm{HH}}(A/R_N,M)
 =\operatorname{Hom}_{R_N}(A^{\otimes_{R_N}n},M).
\]
Their differential is
\begin{align*}
 (bF)(a_1,\ldots,a_{n+1})
 ={}&a_1F(a_2,\ldots,a_{n+1})\\
 &+\sum_{i=1}^n(-1)^i
 F(a_1,\ldots,a_i a_{i+1},\ldots,a_{n+1})\\
 &+(-1)^{n+1}F(a_1,\ldots,a_n)a_{n+1}.
\end{align*}
The coefficient ring is held fixed in every input.

For $e\geq d$, abbreviate $A_{d,N}=U_N(\mathfrak h_d)$.
Give $A_{d,N}$ its $A_{e,N}$-bimodule structure through the quotient.
On
\[
 P^n_{e,d,N}=
 \operatorname{Hom}_{\mathbb C}
  (\Lambda^n\mathfrak h_e,\operatorname{Sym}\mathfrak h_d)
       \otimes R_N,
\]
the coefficient action factors through $\mathfrak h_d$.
Define
\begin{align*}
 (d_{\mathrm{CE}}\phi)(x_0,\ldots,x_n)
 ={}&\sum_i(-1)^i x_i\cdot
       \phi(x_0,\ldots,\widehat x_i,\ldots,x_n)\\
 &+\sum_{i<j}(-1)^{i+j}
 \phi(\{x_i,x_j\},x_0,\ldots,\widehat x_i,\ldots,\widehat x_j,
                                                        \ldots,x_n).
\end{align*}
The bracket occupies the first slot in the second sum.

\begin{theorem}\label{mht:thm:rectangular-rees}
Output PBW symmetrization and input antisymmetrization define a
quasi-isomorphism of $R_N$-complexes
\[
 \operatorname{Alt}_{e,d,N}:
 C^\bullet_{\mathrm{HH}}(A_{e,N}/R_N,A_{d,N})
 \longrightarrow (P^\bullet_{e,d,N},\hbar d_{\mathrm{CE}}),
\]
where
\[
 (\operatorname{Alt}_{e,d,N}F)(x_1,\ldots,x_n)
 =\operatorname{sym}_N^{-1}
   \sum_{\sigma\in S_n}\operatorname{sgn}(\sigma)
       F(x_{\sigma(1)},\ldots,x_{\sigma(n)}).
\]
These maps commute with input refinement, output quotient and
parameter reduction.
\end{theorem}

\begin{proof}
Write $U=A_{e,N}$ and
$\mathfrak l_N=R_N\otimes\mathfrak h_e$ with bracket
$[x,y]_N=\hbar\{x,y\}$. Consider the bimodule complex
$K_n=U\otimes_{R_N}\Lambda^n\mathfrak l_N\otimes_{R_N}U$.
For $X=x_1\wedge\cdots\wedge x_n$, its boundary is
\begin{align*}
 \partial(a\otimes X\otimes b)
 ={}&\sum_i(-1)^{i+1}
 \bigl(ax_i\otimes X_{\widehat i}\otimes b
       -a\otimes X_{\widehat i}\otimes x_i b\bigr)\\
 &+\sum_{i<j}(-1)^{i+j}
 a\otimes[x_i,x_j]_N\wedge X_{\widehat i\widehat j}\otimes b.
\end{align*}
Omitted exterior factors retain their order. The augmentation
$K_0\to U$ is multiplication.

Antisymmetrizing the middle factors maps $K$ into the relative bar
resolution. The endpoint bar faces give the first sum.
Adjacent internal faces cancel in pairs except for their
commutators, giving the second sum. This proves the chain identity.
The antisymmetrization is injective because $\mathfrak l_N$ is a
free PBW summand and factorials are invertible.
Thus $\partial^2=0$ follows from the bar identity.

Filter $K_n$ by the two outer PBW degrees plus $n$.
The associated graded is the diagonal Koszul resolution of
$\operatorname{Sym}_{R_N}\mathfrak l_N$. Its diagonal differences
are independent polynomial variables after a linear change of
coordinates. They form a regular sequence even over the ring
$R_N$ with nilpotents.

A cycle has finite filtration degree.
Lift a primitive of its leading symbol and subtract its boundary.
Induction lowers the degree and proves exactness.

Both $K$ and the relative bar are free bimodule resolutions.
Their comparison lifts the identity and is a chain homotopy
equivalence, by lifting successively through the surjective maps
onto cycles. Applying bimodule Hom into $A_{d,N}$ therefore
preserves that equivalence.

The resulting coefficient action is the commutator with the image
of $x$. For a symmetric polynomial $f$, PBW gives
\[
 [x,\operatorname{sym}_N(f)]
      =\hbar\operatorname{sym}_N(x\cdot f).
\]
The internal bracket term has the same factor $\hbar$.
Hence the differential is precisely $\hbar d_{\mathrm{CE}}$.
All formulas use only the quotient Lie maps and the coefficient
ring, so the three stated naturalities hold.
\end{proof}

The factor $\hbar$ is necessary. For the quadratic jet,
$[H,E]=2\hbar E$, whereas $(d_{\mathrm{CE}}E)(H)=2E$.
The Hochschild differential of $E$ becomes zero on the commutative
special fiber. An unscaled Chevalley differential would not.

The comparison is relative to $R_N$, and its formal limit is
relative to $R=\mathbb C[[\hbar]]$. Already for $\mathfrak l=0$,
the absolute continuous derivation $\partial_\hbar$ on $R$
has no relative counterpart.

\section{Continuous input and output cutoffs}

Form the complete algebra and its quotient tensor powers
\[
 \mathcal A=\varprojlim_{d,N}A_{d,N},\qquad
 \mathcal A^{\widehat\otimes_R n}
      =\varprojlim_{d,N}A_{d,N}^{\otimes_{R_N}n}.
\]
The quotients are discrete, and no PBW-degree completion is taken
at a fixed $(d,N)$. For $n=0$ the tensor source is $R$.
Distinct indices in the tensor factors give the same limit,
because a common upper bound is cofinal.

\begin{definition}\label{mht:def:uniform-cochains}
The continuous cochain space is
\[
 C^n_{\mathrm{uc}}(\mathcal A/R,\mathcal A)
 =\varprojlim_{d,N}\varinjlim_{e\geq d}
  \operatorname{Hom}_{R_N}
   (A_{e,N}^{\otimes_{R_N}n},A_{d,N}).
\]
Input refinement is precomposition with a quotient.
The inverse maps reduce the output and parameter precision.
Each direct limit is discrete, and the cochain space has their
inverse-limit topology.
\end{definition}

\begin{proposition}\label{mht:prop:continuous-cochain-carrier}
There is a natural identification
\[
 C^n_{\mathrm{uc}}(\mathcal A/R,\mathcal A)
 =\operatorname{Hom}_{R,\mathrm{cont}}
    (\mathcal A^{\widehat\otimes_R n},\mathcal A).
\]
All quotient projections of the displayed tensor source are onto.
\end{proposition}

\begin{proof}
Compatible PBW bases identify the underlying module with
$(\varprojlim_d\operatorname{Sym}\mathfrak h_d)[[\hbar]]$.
The tensor powers have the analogous coefficientwise description.
Lift a finite PBW polynomial by the same monomials at all higher
cutoffs. This proves surjectivity. Coefficientwise division also
gives
$\mathcal A/\hbar^{N+1}\mathcal A=\varprojlim_dA_{d,N}$,
and the same identity for each completed tensor.

A continuous map to the discrete output $(d,N)$ factors through
an input quotient $(e,M)$. Enlarge these indices so that
$e\geq d$ and $M\geq N$. The factored map is $R$-linear by
surjectivity of the input projection. Its target is annihilated
by $\hbar^{N+1}$, so the input precision reduces to $N$.

Two factorizations agree after a common refinement.
This gives the stated inverse and direct limits. Conversely,
such a family defines a continuous map to every output quotient.
\end{proof}

The larger input cutoff is essential. Choose a basis Hamiltonian
$x$ of degree $e>d$. Extract its coefficient as a PBW monomial
and multiply by the unit of $\mathcal A$. This is a continuous
linear cochain with scalar output. It detects $x$ and therefore
cannot factor through $A_{d,N}$.

For an arity-$n$ cochain use reduced degree $n-1$.
The brace $F\{G_1,\ldots,G_r\}$ sums over ordered, disjoint
insertions into $F$. If $G_j$ has arity $q_j$ and its block begins
after $i_j$ input arguments, the sign is
$(-1)^{\sum_j(q_j-1)i_j}$. The cup product is
\[
 (F\smile G)(a_1,\ldots,a_{p+q})
   =F(a_1,\ldots,a_p)G(a_{p+1},\ldots,a_{p+q}).
\]

\begin{proposition}\label{mht:prop:continuous-braces}
The Hochschild differential, cup product and every finite brace
are defined on $C^\bullet_{\mathrm{uc}}$ and satisfy their
usual identities.
\end{proposition}

\begin{proof}
Fix output $(d,N)$ and an input cutoff $e$ for the outer cochain.
Each inner cochain is needed only at output $(e,N)$.
Choose a common input cutoff for this finite family.
Every insertion then factors through one tensor quotient, and
further refinement leaves it unchanged.

Multiplication descends at each quotient, so the same argument
defines the differential and cup product. In each identity choose
one common quotient for all intermediate outputs.
The usual insertion bijection and its reduced-degree signs prove
the identity there. Equality at every output proves it in the
complete complex.
\end{proof}

A single space
$\operatorname{Hom}(A_{e,N}^{\otimes n},A_{d,N})$ with $e\ne d$
has no self-valued brace composition. The complete family supplies
the intermediate output required by the proof.

There is a relative derived interpretation with the same topology.
Use complete separated linearly topologized $R$-modules and their
quotient-completed tensors. A complete $\mathcal A$-bimodule has
a continuous action on that tensor. Declare a short exact sequence
admissible when it has a continuous $R$-linear splitting.
This splitting need not commute with the bimodule actions.

\begin{proposition}\label{mht:prop:continuous-relative-bar}
In this exact structure the augmented completed bar is a relative
projective resolution of $\mathcal A$. Consequently,
\[
 \operatorname{RHom}^{\mathrm{rel,cont}}_{\mathcal A^e}
                (\mathcal A,\mathcal A)
 \simeq C^\bullet_{\mathrm{uc}}(\mathcal A/R,\mathcal A),
 \qquad
 \mathcal A^e=\mathcal A\widehat\otimes_R\mathcal A^{\mathrm{op}}.
\]
\end{proposition}

\begin{proof}
Induction from a complete coefficient module $E$ has the
adjunction
\[
 \operatorname{Hom}_{\mathcal A^e,\mathrm{cont}}
       (\mathcal A^e\widehat\otimes_RE,M)
 =\operatorname{Hom}_{R,\mathrm{cont}}(E,M).
\]
The extension sends $a\otimes e\otimes b$ to $a f(e)b$.
Refine first the quotient for the action and then that for $f$.
This proves continuity. Restriction to $1\otimes E\otimes1$
is the inverse. Composing with an admissible splitting proves
relative projectivity of induced modules.

The degree-$n$ bar term is induced from
$\mathcal A^{\widehat\otimes_Rn}$. Its boundary is the alternating
sum of adjacent multiplications. Inserting the unit on the left
gives a continuous $R$-linear contraction of the augmented complex.
On cycles it splits the preceding boundary map, so exactness is
admissible. Continuous bimodule Hom into $\mathcal A$ gives the
cochain spaces and differential above.
\end{proof}

This relative exact structure is part of the assertion.
An absolute topological derived category, or a different analytic
tensor product, requires a separate comparison.

\begin{theorem}\label{mht:thm:continuous-rees}
The finite antisymmetrizations induce a quasi-isomorphism
\[
 C^\bullet_{\mathrm{uc}}(\mathcal A/R,\mathcal A)
 \longrightarrow
 \varprojlim_{d,N}\varinjlim_{e\geq d}
             (P^\bullet_{e,d,N},\hbar d_{\mathrm{CE}}).
\]
For both sides the ordinary inverse limit computes the derived
inverse limit of the indicated direct-limit complexes.
\end{theorem}

\begin{proof}
Filtered direct limits of modules are exact, so
Theorem~\ref{mht:thm:rectangular-rees} gives a quasi-isomorphism
at each output $(d,N)$. Use a countable cofinal diagonal for
these two indices.

Each transition is degreewise surjective. Represent a cochain at
one input cutoff and increase it if necessary. Lift its values
on the PBW tensor basis to the next output and parameter precision.
Extend over the new coefficient ring. For the Chevalley side,
use the exterior basis instead. These are lifts of elements,
not chain sections or $R_N$-linear splittings of coefficient
reduction.

For a degreewise surjective tower, the map
$1-\mathrm{shift}$ on its product is degreewise onto.
Indeed, solve its equations successively using the surjective
transition maps. Its kernel is the ordinary inverse limit, so
the derived-limit fiber complex is quasi-isomorphic to that
kernel. Derived limits preserve the levelwise quasi-isomorphisms,
which proves the comparison.
\end{proof}

At a fixed jet cutoff this reduces to
\[
 C^\bullet_{\mathrm{HH,cont}}
       (U_\hbar(\mathfrak h_d)/R,U_\hbar(\mathfrak h_d))
 \simeq
 \bigl(C^\bullet_{\mathrm{CE}}
   (\mathfrak h_d,\operatorname{Sym}\mathfrak h_d)[[\hbar]],
                                             \hbar d_{\mathrm{CE}}\bigr).
\]
The coefficientwise PBW description is used throughout.
The inverse limit over jets has the additional input cutoff.
A diagonal tower of self-valued finite cochains would omit it.

\section{Integral torsion and even traces}

Scaling a differential changes integral cohomology.
Let $(P^\bullet,d)$ be a complex of complex vector spaces and
write $Z^n=\ker d$, $B^n=\operatorname{im}d$.
All formal series below mean coefficientwise products.

\begin{proposition}\label{mht:prop:scaled-ce-torsion}
There is a natural exact sequence
\[
 0\longrightarrow B^n
 \longrightarrow H^n(P^\bullet[[\hbar]],\hbar d)
 \longrightarrow H^n(P^\bullet,d)[[\hbar]]
 \longrightarrow0.
\]
The first term is annihilated by $\hbar$ and is exactly the
$\hbar$-power torsion. A vector-space complement of $B^n$ in
$Z^n$ splits the sequence. The first Bockstein of the
special-fiber complex, whose differential is zero, is $d$.
\end{proposition}

\begin{proof}
Multiplication by $\hbar$ is injective on every coefficient
module. Cycles and boundaries are therefore
$Z^n[[\hbar]]$ and $\hbar B^n[[\hbar]]$.
Coefficientwise passage to $Z^n/B^n$ gives the surjection, with
kernel $B^n[[\hbar]]/\hbar B^n[[\hbar]]=B^n$.
A complement lifts the coefficients and splits it.

The quotient has no $\hbar$-power torsion, while the kernel is
killed by $\hbar$. Finally lift $x\in P^n$ by its constant series.
Its differential is $\hbar dx$. Dividing by $\hbar$ and reducing
gives the Bockstein $dx$.
\end{proof}

For example, $R\xrightarrow{\hbar}R$ has free cochain modules and
first cohomology $R/(\hbar)$. At finite precision its degree-zero
kernel is $\mathbb C\hbar^N$. The next reduction sends that
kernel to zero. Those extra cycles disappear in the limit,
while the degree-one torsion class persists.

Unimodularity controls the first adjoint trace, but the even traces
have another limit. Fix $m\geq1$. For $x\in\mathfrak h$, let $x_2$ be its
quadratic part and let $\rho_r$ denote its action on $W_r$.
Higher Hamiltonians raise degree, so
\[
 T_{2m,d}(x)=\operatorname{Tr}_{\mathfrak h_d}
       (\operatorname{ad}_x^{2m})
    =\sum_{r=2}^d\operatorname{Tr}_{W_r}(\rho_r(x_2)^{2m}).
\]

\begin{proposition}\label{mht:prop:hamiltonian-even-traces}
The expression $T_{2m,d}$ is polynomial in the cutoff $d$.
Its constant coefficient, denoted $\operatorname{FP}_d$, is
\[
 \operatorname{FP}_dT_{2m,d}(x)
      =-\operatorname{Tr}_{W_1}(\rho_1(x_2)^{2m}).
\]
For $H=-uv$,
\[
 T_{2,d}(H)=\frac{d(d+1)(d+2)(d+3)}{12}-2.
\]
\end{proposition}

\begin{proof}
For $x_2=aH$ the sum is
$a^{2m}\sum_{r=2}^d\sum_{j=0}^r(r-2j)^{2m}$.
Power-sum identities make it polynomial in $d$.
The sum starting at $r=1$ vanishes at $d=0$. Removing its first
term subtracts $2a^{2m}$. This proves the constant coefficient
on the Cartan line. Both sides are invariant polynomials on
$\mathfrak{sl}_2$, so density of semisimple elements proves
the identity everywhere.

For $m=1$ the inner sum is $r(r+1)(r+2)/3$.
Summing from $r=2$ gives the displayed polynomial.
\end{proof}

This defines a subtraction scheme, not a canonical limit.
Let $B_{2m}$ be the Bernoulli numbers defined by
$z/(e^z-1)=\sum_n B_nz^n/n!$. The formal half-determinant
logarithm has coefficients
$c_m=B_{2m}/(4m(2m)!)$. Subtracting the nonconstant cutoff
polynomial from each trace gives
\[
 \ell_{\mathrm{ren}}(x)
 =-\frac12\operatorname{Tr}_{W_1}
  \log\!\left(
   \frac{\sinh(\hbar\rho_1(x_2)/2)}{\hbar\rho_1(x_2)/2}
       \right).
\]
The quotient is its power series with constant term one, without
inverting an eigenvalue. On the Cartan line its exponential is
$(\hbar a/2)/\sinh(\hbar a/2)$.

Replacing the cutoff by $d+b$ changes the two-wheel finite part
by $b(b+1)(b+2)(b+3)/12$.

An ordinary even dual has the same even traces and doubles them
when added to the original representation. A parity-reversed
dual cancels them in the supertrace. These two coefficient
choices therefore have different subtraction problems.

\section{Support and the remaining field comparison}

Let $\mathcal L_c(U)$ be compactly supported smooth sections of a
local differential graded Lie algebra on an open set $U$.
Assume its differential and bracket are differential operators.
Extension by zero then gives a covariant Lie map
$\mathcal L_c(U)\to\mathcal L_c(V)$ for $U\subset V$.
For disjoint opens, the bracket of their extended sections is zero.

With the algebraic tensor product, the Chevalley chains
$C^{\mathrm{Lie}}_*(\mathfrak l)=\operatorname{Sym}(\mathfrak l[1])$
have their internal and bracket differential. They take direct
sums of Lie algebras to tensor products. Consequently the maps
\[
 \bigotimes_i C^{\mathrm{Lie}}_*(\mathcal L_c(U_i))
 \longrightarrow C^{\mathrm{Lie}}_*(\mathcal L_c(V))
\]
define a prefactorization algebra. Weiss descent is an additional
property.

Extension by zero pulls Chevalley cochains from $V$ to $U$.
A covariant observable model therefore requires its own coefficient
duality and support construction.

For the usual smooth topologies,
$(C^\infty(U))'_b=\mathcal E'(U)$ consists of compactly supported
distributions. In contrast,
$(C^\infty_c(U))'_b=\mathcal D'(U)$ allows arbitrary support.
The first assertion follows because a continuous functional uses
finitely many seminorms on a common compact set.
The second follows by compatibility on all compact test supports.
The subscript indicates the strong dual topology.

The Hochschild side of Theorem~\ref{mht:thm:continuous-rees}
carries braces and a relative derived-center interpretation.
An operation-preserving comparison with independently defined
polyvectors requires further maps.

\begin{problem}\label{mht:qn:continuous-hamiltonian-comparison}
Construct a coherent refinement of the continuous Rees comparison
for independently defined polyvector products and brackets.
Require compatibility with all input and output refinements,
and specify any even-wheel subtraction used by a multiplicative
PBW correction.

Construct a map from supported mixed field observables to this
coefficient complex. Preserve horizontal jets, normal derivatives,
the vacuum and collision operations. Make the subtraction compatible
with these operations and nonlinear symplectic coordinate descent.
\end{problem}
